\section{Scope-Bounding and Damage Guarantees}
\label{sec:damage_guarantees}

Section~\ref{sec:relaxation_protocol} treated the damage tolerance
$\damagebound$ as a parameter; here we show that the scope
construction makes $\damagebound$ an enforceable upper bound on
the worst-case $\volL$-contraction during and after a test, even
when cooperative capabilities can propagate contraction outside
the nominal scope.

\subsection{Formalizing test scope}

\begin{definition}[Test Scope]
\label{def:test_scope}
A \emph{test scope} $\scope$ for risk claim
$R = (a_j, (d_1, \ldots, d_k), \Pathway, \pR)$ is a sub-poset of
$P$ satisfying:
\begin{itemize}
\item[(i)] \textbf{Pathway containment.}
  $\Pathway$ is exercisable within $\scope$: all prerequisite
  capabilities are present in $\scope$, and all intermediate
  poset transitions along $\Pathway$ remain within $\scope$.
\item[(ii)] \textbf{Forward-reachable cooperative closure.}
  For every cooperative capability $c$ in the full poset $P$ that
  involves at least one capability in $\scope$ and at least one
  capability in $P \setminus \scope$, the maximum
  $\volL$-contraction along any pathway within $\scope$ that
  modifies $c$ is included in $\maxcon(\scope)$.  Furthermore, if
  modification of $c$ can causally enable (unlock) an exercise
  pathway in $P \setminus \scope$ (that is, there exists a
  pathway $\Pathway_{\mathrm{ext}}$ in $P \setminus \scope$ whose
  prerequisites include a state change in $c$), then the expected
  contraction from $\Pathway_{\mathrm{ext}}$, weighted by its
  activation probability, is also included in $\maxcon(\scope)$.
  This transitive-closure condition prevents a test within
  $\scope$ from triggering a cascade outside $\scope$ through
  cross-boundary cooperative capabilities.
\item[(iii)] \textbf{Irreversibility bound.}
  For each exercise pathway within $\scope$, the fraction of the
  pathway's $\volL$-contraction that is irreversible (cannot be
  recovered by reinstating $\Xirestrict$) is bounded by a known
  constant $\deltairrev(\Pathway) \leq \damagebound$.  Pathways
  with $\deltairrev > \damagebound$ are excluded from the scope.
\item[(iv)] \textbf{Scope minimality.}
  $\scope$ is inclusion-minimal subject to
  (i)--(iii): no proper sub-poset of $\scope$ satisfies all three
  conditions.
\end{itemize}
\end{definition}

\begin{remark}[Constructibility of the scope]
\label{rem:constructibility}
Condition~(i) is computable from the exercise-pathway definition
of~\citep[Definition~2]{lasser2026horizon}: enumerate the
prerequisite set and all intermediate states.  Condition~(ii)
requires computing the cooperative capabilities that span
$\scope$ and $P \setminus \scope$ (the \emph{cooperative
surface} of $\scope$) and then forward-reachability analysis:
for each boundary cooperative, enumerate external pathways whose
prerequisites include a state change in that cooperative.  This
is more expensive than simple pathway enumeration (exponential
in the depth of the forward-reachability chain) but is necessary
to prevent cascading damage through cooperative-capability
spillover.  In practice the forward-reachability depth is
bounded by the poset diameter and most cooperative capabilities
have low fan-out.  Condition~(iii) requires a pathway-level
reversibility analysis; most capability restrictions are
reversible (reinstatement restores the prior state), but
information-disclosure and trust-degradation pathways may have
irreversible components.
\end{remark}

\subsection{The damage-bound theorem}

\begin{theorem}[Damage Bound for Controlled Relaxation]
\label{thm:damage_bound}
Let $T = (R, \Xirestrict, \scope, \testdur, \damagebound)$ be a
controlled relaxation test with test scope $\scope$ satisfying
Definition~\ref{def:test_scope}.  Then:
\begin{itemize}
\item[(a)] \textbf{Maximum instantaneous damage.}  At any step
  $t' \in \{t, t+1, \ldots, t + \testdur\}$ during the test,
  $$
  \volL(G_{t'}) \geq \volL(G_t) - \maxcon(\scope),
  $$
  where $\maxcon(\scope)$ is the maximum $\volL$-contraction
  achievable by any exercise pathway within $\scope$.  Under the
  scope-selection constraint
  $\maxcon(\scope) \leq \damagebound$, the instantaneous damage
  is bounded by $\damagebound$.
\item[(b)] \textbf{Partial reversibility under reinstatement.}
  If the safety cutoff triggers at some step
  $t' < t + \testdur$ (the detector reports contraction within
  $\scope$ exceeding $\damagebound$; see
  Section~\ref{sec:relaxation_protocol} Phase~2 for the
  single-detector convention), reinstating $\Xirestrict$ at
  $t' + 1$ gives
  $$
  \volL(G_{t'+1}) \geq \volL(G_t) - \damagebound
       - \deltairrev^{\max},
  $$
  where
  $\deltairrev^{\max} = \max_{\Pathway' \subseteq \scope}
      \deltairrev(\Pathway')$
  is the maximum irreversible component across all pathways in
  the scope (Definition~\ref{def:test_scope}(iii)).  By the scope
  construction $\deltairrev^{\max} \leq \damagebound$, so total
  post-reinstatement damage is bounded by $2\damagebound$.
\item[(c)] \textbf{Scope isolation under cooperative closure.}
  If no other test is running on an overlapping scope
  (Definition~\ref{def:cr_test}, Phase~1, condition~(iii))
  and $\scope$ satisfies cooperative closure
  (Definition~\ref{def:test_scope}(ii)), then a cooperative-closed
  $\scope$ exists (Lemma~\ref{lem:computable_maxcon} gives a
  constructive enumeration of the cooperative surface and its
  forward-reachability closure), the cross-boundary contribution
  to contraction is finite and enumerable, and the damage from
  test~$T$ that propagates to $P \setminus \scope$ is bounded by
  the cooperative-surface terms that Definition~\ref{def:test_scope}(ii)
  requires to be included in $\maxcon(\scope)$.  Formally: for
  any capability $d \notin \scope$ that is not part of any
  cooperative capability spanning $\scope$ and $P \setminus \scope$,
  the exercise probabilities of $d$'s pathways are unchanged
  during the test.  The substantive content of part~(c) is
  therefore the existence and tractability of the closure; the
  containment bound follows by construction from
  Definition~\ref{def:test_scope}(ii).
\end{itemize}
\end{theorem}

\begin{proof}[Proof sketch]
Part~(a) follows from the pathway-enumeration construction of
$\maxcon$ under finite pathway length, and the scope-selection
constraint.  Part~(b) decomposes the contraction at the cutoff
step into its reversible and irreversible components; the
reversible component is restored by reinstating $\Xirestrict$,
the irreversible component is bounded above by $\damagebound$
under Definition~\ref{def:test_scope}(iii).  Part~(c) is the
direct consequence of cooperative closure: cross-boundary
cooperative terms are accounted for in $\maxcon(\scope)$ by
construction, and all other boundary capabilities are untouched
by the test.  Full derivation is in Appendix~\ref{app:proofs},
Section~\ref{app:proof_damage_bound}.
\end{proof}

\begin{remark}[Scope existence as a scheduling precondition]
\label{rem:scope_existence}
Theorem~\ref{thm:damage_bound} is stated conditionally on the
existence of a scope satisfying Definition~\ref{def:test_scope}.
For posets with dense cross-boundary cooperative structure, the
cooperative-closure requirement may force $\scope = P$, in which
case $\maxcon(P) \leq \damagebound$ may fail and no test is
schedulable.  The feasibility condition is enforced by the
meta-policy (Section~\ref{sec:meta_policy}, Definition~6,
Phase~1); Theorem~\ref{thm:damage_bound} applies only when the
meta-policy reports a feasible scope.  When no feasible scope
exists, the restriction persists and the Wamura problem is not
resolved for that specific risk claim by controlled relaxation;
alternative mechanisms (structural discovery via
\citep[Proposition~2]{lasser2026horizon}, or scope-refinement
that trades coverage for feasibility) must apply.
\end{remark}

\begin{remark}[Why not $\deltairrev \leq 0$?]
\label{rem:why_not_zero}
The preemptive restriction
criterion~\citep[Proposition~1]{lasser2026horizon} is an ex-ante
cost-benefit comparison justifying the initial restriction; it
does not guarantee that reinstatement after partial pathway
execution fully recovers the pre-test state.  Irreversible state
changes during the test (information disclosure to adversaries,
trust degradation, poset structural changes not reversed by
restriction alone) create a gap between the pre-test $\volL$ and
the post-reinstatement $\volL$.
Definition~\ref{def:test_scope}(iii) bounds this gap by
requiring the scope to exclude pathways with large irreversible
components.
\end{remark}

\begin{lemma}[Computable max\_contraction]
\label{lem:computable_maxcon}
For a finite sub-poset $\scope$ with bounded pathway length $L$
and bounded cooperative-surface forward-reachability depth
$L'$, $\maxcon(\scope)$ is computable in
$O(|\scope|^{L+1} + |P|^{L'})$ time by exhaustive pathway
enumeration together with enumeration of cross-boundary
cooperative chains.  When the simplifying assumption $L' \leq L$
holds (cross-boundary cascades cannot chain more than
pathway-length-many cooperative hops), the bound simplifies to
$O(|P|^{L+1})$.
\end{lemma}

\begin{proof}[Proof sketch]
Pathway length is bounded
in~\citep[Definition~2]{lasser2026horizon}.  The number of
internal pathways is at most $|\scope|^L$ (each step selects a
capability in $\scope$); each pathway contraction is computable
from the axioms M1--M6~\citep[Proposition~1]{lasser2026poset} in
$O(|\scope|)$ time.  The cooperative-surface terms from
Definition~\ref{def:test_scope}(ii) require enumerating
cross-boundary cooperative chains to depth $L'$, costing
$O(|P|^{L'})$ in the worst case.  Take the maximum across
internal pathways and add cooperative-surface terms.
\end{proof}

\begin{remark}[Computational tractability]
\label{rem:computability}
The bound is exponential in pathway length and
forward-reachability depth but polynomial in scope size for
fixed $L, L'$.  Practical deployment
requires either (a)~short pathway restrictions ($L \leq 3$,
say), (b)~approximate bounds using the concentration-risk
formula
of~\citep[Definition~3]{lasser2026horizon} as an upper bound on
$\maxcon$, or (c)~structural decomposition of $\scope$ into
independent sub-scopes.  This is an engineering consideration,
not a theoretical limitation; the bound exists and is
computable; the question is computational cost.
\end{remark}

\begin{corollary}[Damage bound under conservative over-approximation]
\label{cor:over_approx_safety}
Let $\maxcon^\star(\scope) \geq \maxcon(\scope)$ be any
over-approximation of the maximum scope contraction.  A
scope-selection rule that requires $\maxcon^\star(\scope) \leq
\damagebound$ continues to satisfy the damage bound of
Theorem~\ref{thm:damage_bound}: worst-case $\volL$-contraction
during the test is $\leq \maxcon(\scope) \leq
\maxcon^\star(\scope) \leq \damagebound$, and the post-cutoff
bound $2\damagebound$ inherits unchanged.  Cooperative-closure
enumeration cost is therefore tradeable against scope
conservatism: deployers may accept a narrower $\scope$ chosen
under a cheaper-to-compute over-approximation of the
cooperative surface and still preserve the safety guarantee,
at the cost of feasibility (some otherwise-feasible scopes will
be rejected).
\end{corollary}

\begin{proof}[Proof sketch]
Direct.  By assumption $\maxcon(\scope) \leq
\maxcon^\star(\scope) \leq \damagebound$, so
Theorem~\ref{thm:damage_bound}(a)'s inequality $\volL(G_{t'})
\geq \volL(G_t) - \maxcon(\scope) \geq \volL(G_t) -
\damagebound$ holds.  Part~(b) inherits because
the irreversibility bound
$\deltairrev^{\max} \leq \damagebound$ of
Definition~\ref{def:test_scope}(iii) is a pathway-level
condition unaffected by the maxcon over-approximation.
Part~(c) inherits because cooperative-surface contractions are
included inside $\maxcon$ by construction
(Definition~\ref{def:test_scope}(ii)) and
$\maxcon \leq \maxcon^\star \leq \damagebound$.
\end{proof}
